\section{Introduction}\label{sec:intro}

An agent that operates a website has to perceive it, and the two standard ways are complementary in principle: the page's text (an accessibility tree) is precise and cheap, a screenshot shows what the text omits. A third option fuses them, overlaying numbered marks on the screenshot and listing the marked elements as text \citep{yang2023som,koh2024visualwebarena}. Practitioners typically pick one of the three for a whole deployment. The question this paper asks is the one that choice raises: \emph{is the choice worth making per task, and can it be made?}

Answering it requires keeping apart three quantities that a single leaderboard number merges:
\begin{enumerate}\setlength{\itemsep}{0pt}
  \item \textbf{Representation value}: does the screenshot change success, holding everything else fixed?
  \item \textbf{Preference reliability}: when one representation beats another on a task, is that a stable property of the task, or a property of one run?
  \item \textbf{Routing gain}: can a policy that chooses per task, using only what is known before it runs, beat the fixed choices?
\end{enumerate}
A hindsight ceiling picks, for each task, whichever representation happened to succeed. However large, it answers none of them, because it rewards run-to-run luck as much as structure. We therefore rerun identical conditions and read every comparison against what a rerun alone changes.

\paragraph{Findings.}
(i) No choice wins everywhere: the fused mode leads on classifieds, text leads on WebArena reddit, and the order reverses across backbones (\S\ref{sec:deploy}).
(ii) The screenshot has a specific value that can be isolated and predicted: with the text payload and prompt held fixed, adding the screenshot raises success on the classifieds tasks a zero-cost intent rule flags by 25--30 points for the two strongest backbones (11 for a 4B model), against 5--10 on the other tasks; the effect reproduces on a full rerun and is larger for the stronger backbones we tested. On reddit the same rule finds nothing (\S\ref{sec:screenshot}).
(iii) Per-task preferences between representations are reproducible beyond an additive difficulty model: a no-interaction null conditioned on every task's and every mode's success count is rejected in every fully rerun cell. Yet a single run measures them with reliability 0.31--0.50, and a decision study projects that labels averaged over two to three runs would pass 0.5 (\S\ref{sec:preference}).
(iv) Two lightweight routers that fix their operating point on training data show no reliable gain over the fixed representations and their random mixtures; the one hindsight opportunity that passes its null comes from a router that only decides which tasks deserve the expensive mode, holds under billed cost only, and is not captured without looking at test outcomes (\S\ref{sec:routing}).

\paragraph{Contributions.}
A rerun-controlled measurement of text, screenshot and fused observation across three sites and four backbones, with text-only ablations that isolate the screenshot; a matched estimate of the screenshot's value and an ex-ante rule that predicts it on one site; a test of whether per-task representation preferences are reproducible, with a decision study of how many runs a reliable label needs; and a deployable-policy comparison against fixed choices \emph{and their random mixtures}, which is the comparison a routing claim has to win; we take this baseline from model-routing benchmarks and apply it to observation choice (\S\ref{sec:method}).
